\section{Висновки}
\label{sec:conclusions}

Побудовано аналітичну модель ініціювання відгалуження від заданої
стаціонарної V-подібної міжфазної тріщини зсуву в кутовій точці ламаної
межі поділу двох пружних ізотропних матеріалів. Дві симетричні гілки
тріщини мають контактуючі без тертя береги, а рівень зовнішнього
навантаження вважається недостатнім для втрати локальної стійкості біля
їхніх зовнішніх вершин. Локальне руйнування біля кутової точки описано
малою розтягувальною зоною передруйнування, орієнтованою вздовж
бісектриси кутового сектора одного з матеріалів і змодельованою лінією
розриву нормального переміщення зі сталим відривним напруженням. За допомогою перетворення
Мелліна, зрощування асимптотик і методу Вінера--Гопфа отримано замкнені
співвідношення для довжини зони $d_i$, її розкриття $\delta_i$ та
зонального енергетичного параметра $J_i$.

Фізична допустимість цієї локальної конфігурації визначається
одночасним виконанням двох умов,
\begin{equation*}
  Cg_2<0,\qquad CQ_i>0.
\end{equation*}
Перша з них забезпечує стискальний контакт берегів початкової міжфазної
тріщини, а друга --- розтягувальне нормальне поле в матеріалі, де може
утворитися зона передруйнування. Для базової пари
$E_1/E_2=0.5$, $\nu_1=\nu_2=0.3$ ці умови розбивають інтервал кутів на
фізично допустимі гілки з переходами
$\alpha_1=12.9202^\circ$ і $\alpha_2=107.0897^\circ$ та однозначно
визначають, у якому з двох матеріалів може виникнути зона за заданого
знака $C$. Розрахунки для інших значень $E_1/E_2<1$ показують, що
положення цих переходів закономірно змінюється з пружним контрастом
матеріалів.

Аналітичні формули встановлюють степеневу залежність розміру зони від
зовнішнього навантаження та опору матеріалу відриву. На кожній фізично
допустимій гілці зі збільшенням модуля нормованого навантаження
$|\sigma'|$ величини $d_i/l$, $\delta_i'$ і $J_i'$ зростають узгоджено;
при фіксованій геометрії $\delta_i$ і $J_i$ пропорційні довжині зони.
Зменшення опору відриву $\sigma_i$ збільшує як довжину, так і розкриття
зони. На широких фізично допустимих гілках кутові залежності трьох
параметрів мають близькі, але не тотожні максимуми, оскільки коефіцієнти
пропорційності між $d_i/l$, $\delta_i'$ та $J_i'$ самі залежать від
кута. Ці максимуми характеризують геометрії, для яких за заданого
зовнішнього навантаження досягаються найбільші локальні деформаційні та
енергетичні параметри. Для визначення кутової залежності критичного
навантаження необхідно окремо задати відповідний критичний матеріальний
параметр і для кожного кута розв'язати рівняння критерію.

Порівняння інтактної кутової точки, конфігурації із заданою
міжфазною зсувною тріщиною та конфігурації тріщина--зона
передруйнування показує зміну порядку кутової сингулярності. Для
розглянутої базової пари матеріалів на фізично допустимих сегментах
конфігурація з тріщиною має сильнішу концентрацію напружень, тоді як
додавання зони передруйнування частково її екранує; при цьому виконується
\begin{equation*}
  |\lambda_0|<|\lambda_i|<|\lambda|.
\end{equation*}
Отже, сама сингулярність створює передумови для локального руйнування,
але не є достатнім критерієм відгалуження: необхідно також враховувати
умови допустимості, опір відриву та рівень зовнішнього навантаження.

У межах прийнятої локальної моделі величину $J_i$ використовуємо як
зональний енергетичний параметр. Після незалежного задання критичної
характеристики матеріалу $J_{i\mathrm{c}}$ умова
$J_i=J_{i\mathrm{c}}$ визначає критичне зовнішнє навантаження
ініціювання відгалуження нормального відриву (мода~I). Альтернативно
може бути використано
деформаційний критерій $\delta_i=\delta_{i\mathrm{c}}$; за фіксованої
геометрії обидва критерії описують ту саму однопараметричну сім'ю
рівноважних станів і є еквівалентними лише за узгодженого вибору
критичних матеріальних параметрів.

Кількісні результати мають інтерпретуватися в межах припущень моделі:
$d_i\ll l$, плоска деформація, симетричне навантаження, відсутність
тертя та використання провідного сингулярного члена зовнішнього поля.
У частині кутового діапазону розраховане значення $d_i/l$ досягає
приблизно $0.18$, тому поблизу відповідних максимумів маломасштабне
припущення стає граничним і кількісні оцінки потребують обережності.
Модель описує лише ініціювання відгалуження; подальший скінченний шлях і
стійке поширення вторинної тріщини цією постановкою не охоплюються.
